\documentclass[11pt,a4paper]{article}

\usepackage[margin=1.05in]{geometry}
\usepackage[T1]{fontenc}
\usepackage{charter}
\usepackage[scaled=0.9]{helvet}
\usepackage{courier}

\usepackage{amsmath,amssymb}
\usepackage{booktabs}
\usepackage{array}
\usepackage{tabularx}
\usepackage{longtable}
\usepackage{xcolor}
\usepackage{listings}
\usepackage[colorlinks=true,linkcolor=black,citecolor=black,urlcolor=blue]{hyperref}
\usepackage{enumitem}
\usepackage[protrusion=true,expansion=false]{microtype}
\setlength{\emergencystretch}{3em}
\DisableLigatures[-,<,>]{encoding = T1, family = tt*}
\setlist{itemsep=2pt,topsep=4pt}

% NOTATION RULES (as in the note and the spec): the rho calculus is never
% written with the Greek letter; quoting is @P, dropping is *x.
\newcommand{\rhoc}{\textsf{rho}}
\newcommand{\nm}[3]{\langle @#1,\,#2,\,#3\rangle}
\newcommand{\tool}{\texttt{f1r3score}}
\newcommand{\ttx}[1]{\texttt{#1}}

\definecolor{codebg}{gray}{0.96}
\definecolor{notebg}{rgb}{0.97,0.95,0.88}
\lstdefinelanguage{score}{
  morekeywords={score,pitches,durations,timbres,def,factor,machine,table,
    for,where,par,in,at,play,line,then,import,pitchset,base,keyloc,keycode,
    proc,name,timbre,pitch,dur,clause,cont,string,true,false,nu,mu,held,carry,
    prev,step,back,last,call,passes,leads,live,any,all,sum,uniform,dual,scale,
    gm,on,piano88},
  sensitive=true,
  morecomment=[l]{//},
  morecomment=[s]{/*}{*/},
  morestring=[b]",
}
\lstdefinelanguage{Rust}{
  morekeywords={pub,struct,enum,impl,fn,trait,let,mut,use,for,where,self,Self,
    Box,Vec,Result,Ok,Err,Some,None,match,if,else,return,mod},
  sensitive=true, morecomment=[l]{//}, morestring=[b]",
}
\lstset{
  basicstyle=\ttfamily\small,
  backgroundcolor=\color{codebg},
  frame=none,
  keywordstyle=\bfseries,
  commentstyle=\itshape\color{black!60},
  columns=fullflexible,
  keepspaces=true,
  showstringspaces=false,
  xleftmargin=6pt,
  aboveskip=6pt, belowskip=6pt,
  breaklines=true,
  upquote=true,
}
\lstnewenvironment{shell}{\lstset{language=bash,keywordstyle=\ttfamily}}{}
\lstnewenvironment{out}{\lstset{language={},backgroundcolor=\color{white},
  frame=l,framerule=0.6pt,rulecolor=\color{black!35},basicstyle=\ttfamily\footnotesize}}{}

\newenvironment{caveat}[1]{\par\medskip\noindent\colorbox{notebg}{\parbox{\dimexpr\linewidth-2\fboxsep}{\small\textbf{#1.}\ }}\par\nopagebreak\begingroup\small}{\endgroup\par\medskip}

\title{\bfseries F1R3Score\\[4pt]
\large Building, deploying and using \texttt{f1r3score},\\
the player for the score calculus}

\author{F1R3FLY.io}

\date{Documentation --- 27 September 2026\\[2pt]
\normalsize Against \texttt{F1R3FLY-io/F1R3Score} commit \texttt{7416bc0},
which implements revision~2 of \emph{The f1r3score Player} specification}

\begin{document}
\maketitle

\begin{abstract}
\noindent \tool{} is a Rust command-line player and library for the score
calculus of \emph{Scores as Processes}. A score is a process in a variant of
the \rhoc{} calculus, and a performance is one run of it. The player reads a
textual score, elaborates it to a closed core term, and executes a run. The
run can be written out as a table, JSON, a replayable trace of records, a
Standard MIDI File, or live MIDI.

This guide covers three things. It shows how to build the workspace and run its
conformance suite. It covers deploying the player as a binary, as a library, and
as a reproducible archive of performances. And it shows how to use the player
and write scores for it. Every command and every output quoted here was run
against the commit named above, and the guide says where it could not verify
something. Section~\ref{sec:limits} lists where the implementation departs from
the specification or from its own documentation.
\end{abstract}

\tableofcontents
\newpage

%==============================================================================
\section{Introduction}\label{sec:intro}

\subsection{What \tool{} is}

The score calculus proposes a musical notation whose primary objects are
exchange and preference rather than a single fixed performance: who listens to
whom, which continuation a style favours, what lies easily under a player's
hands. \tool{} makes that notation playable. It has three layers:
\begin{itemize}
\item a \emph{surface syntax}, a dialect of the f1r3\textbar@ng control-flow
  sublanguage. It has decorated names, alphabet declarations and an
  elaboration-time macro layer, with files ending in \ttx{.score};
\item a \emph{standard library}, \ttx{std.score}, written entirely in that
  syntax. It defines written notes, reflective voices, hockets, call and
  response, keys, keyboards and chimeric instruments;
\item an \emph{engine} that knows nothing of voices, keys or idioms. It
  matches receipts and messages, weighs the candidates with graded
  where-clauses, resolves contention by maximal matchings, keeps the record of
  every communication, and reads notes off records.
\end{itemize}

\subsection{Sources}

The guide draws on three sources in the \ttx{F1R3FLY-io} GitHub organisation:
\begin{itemize}
\item \ttx{F1R3Score}, the implementation. It contains \ttx{README.md},
  \ttx{DESIGN.md} (decisions, deviations and verification status), the
  examples, and a copy of the note in \ttx{docs/};
\item \ttx{publications/f1r3score/score-calculus.tex}, the note \emph{Scores as
  Processes}, which defines the calculus;
\item \ttx{publications/f1r3score/f1r3score-player-spec.tex}, the player
  specification (revision~2), together with \ttx{skeinsim.py}, the reference
  simulator.
\end{itemize}
Section numbers of the form ``\S5.2 of the note'' refer to \emph{Scores as
Processes}. Requirement and decision numbers refer to the specification.

\subsection{How this guide was checked}

The guide was written against commit \ttx{7416bc0} on Ubuntu~24.04 (x86\_64),
with Rust~1.75.0 from the distribution packages. That is also the workspace's
declared minimum, \ttx{rust-version = "1.75"}. On that machine:
\begin{itemize}
\item the release build succeeds with no warnings, and so does the build with
  the \ttx{live} feature;
\item \ttx{cargo test --release} passes all 70 tests;
\item the cross-implementation differential against \ttx{skeinsim.py} agrees
  on every scenario (run with 10 random scores);
\item every command in Sections~\ref{sec:use} and~\ref{sec:examples} was run,
  and the outputs shown are the actual outputs.
\end{itemize}
Two things could not be verified: the \ttx{wasm32} build of the core crates,
and live MIDI on a real device (\S\ref{sec:limits}).

%==============================================================================
\section{Concepts in brief}\label{sec:concepts}

This section gives the minimum of the calculus needed to read the rest of the
guide. The note is the authority.

\paragraph{Names carry notes.} A name is a triple $\nm{P}{\tau}{d}$: a quoted
process $@P$ (its \emph{location}), a timbre $\tau$, and a datum $d$. The datum
is either a pitch or a duration, and either decoration component may be left
open, written \ttx{\_}. A name whose datum is a pitch is a \emph{channel}; a
name whose datum is a duration is a \emph{co-channel}. The \emph{general name}
$@P$ is $\nm{P}{\_}{\_}$.

\paragraph{Every communication is a note.} A receipt
\ttx{for (y <- x where psi) \{ P \}} and a message \ttx{x'!(Q, tau, d)}
communicate when their subjects \emph{match}. That means three things hold:
their locations are equivalent, their timbres meet in a concrete timbre, and
their data are concrete and of opposite polarity (one a pitch, the other a
duration). Each firing yields a \emph{record}, and the \emph{playback} function
$\nu$ reads a note off it: the timbre, the pitch, and the duration. There are
no silent steps.

\paragraph{Null notes.} The rest \ttx{r} and the null duration \ttx{eps}
(length $0$) let the calculus do bookkeeping, such as a server fetching its
own code, as communications that are notes of zero length. These \emph{null
notes} are recorded but excluded from the performance.

\paragraph{Clauses weigh choices.} A receipt's where-clause evaluates to an
exact non-negative rational on each candidate. A candidate is \emph{available}
when its weight is positive. Available candidates that share a receipt or a
message form a \emph{contention set}. A set is resolved by choosing a maximal
matching with probability proportional to the product of its weights, and
firing all of it. Clauses multiply, so style, register and the physicality of
an instrument can be separate factors whose product is the voice's law.

\paragraph{Time.} Every top-level process carries a stamp. A communication's
onset is the later of its two stamps, and its continuation is stamped at the
onset plus the note's length. Sets are resolved in onset order.

\paragraph{Voices are reflective.} There is no recursion former. An unbounded
voice is a server that keeps its own code on a code channel and re-sends it
each time it is read (\S5.2 of the note). Definitions in the surface syntax are
therefore macros with an acyclic call graph.

\paragraph{Musicians and instruments.} A \emph{key} is a receipt holding a
pitch, and it belongs to the instrument. A \emph{touch} is a message holding a
duration, and it belongs to the player. A keyboard installs keys at key
locations. A player sends touches, and \ttx{;} makes a player wait until its
note has ended (Definition~9.3 of the note). A \emph{chimera} holds a key per
timbre at one location, so a touch that leaves its timbre open is answered in
whichever timbre the keys' clauses favour.

%==============================================================================
\section{Building}\label{sec:build}

\subsection{Prerequisites}

\begin{center}
\small
\begin{tabularx}{\linewidth}{@{}l>{\raggedright\arraybackslash}X@{}}
\toprule
need & for \\
\midrule
Rust $\geq$ 1.75 and Cargo & everything. Any of rustup, the distribution
  packages (Ubuntu 24.04: \ttx{apt install rustc cargo} gives 1.75.0), or
  Homebrew will do. \\
network access to crates.io & the first build only. The default build
  depends only on \ttx{num-bigint}, \ttx{num-rational}, \ttx{num-traits} and
  \ttx{num-integer}. \\
ALSA headers and \ttx{pkg-config} (Linux) & the optional \ttx{live} feature
  only: \ttx{apt install libasound2-dev pkg-config}. On macOS \ttx{midir} uses
  CoreMIDI, and on Windows WinMM; neither needs extra packages. \\
Python 3 and a checkout of \ttx{publications} & the cross-implementation
  differential only (\S\ref{ssec:differential}). \\
a \ttx{wasm32-unknown-unknown} standard library & building the core crates
  for the browser (\S\ref{ssec:wasm}). \\
\bottomrule
\end{tabularx}
\end{center}

\subsection{Getting the source and building}

\begin{shell}
git clone https://github.com/F1R3FLY-io/F1R3Score.git
cd F1R3Score
cargo build --release
./target/release/f1r3score --help
\end{shell}

\noindent A cold release build takes about a minute. The binary is
\ttx{target/release/f1r3score}. It is self-contained: \ttx{std.score} is
compiled in with \ttx{include\_str!}, so no data files need to accompany it.

\subsection{Cargo features}

\begin{center}
\small
\begin{tabularx}{\linewidth}{@{}lll>{\raggedright\arraybackslash}X@{}}
\toprule
feature & crate & default & effect \\
\midrule
\ttx{spigot} & \ttx{score-chance} & on & exact digit streams of
  $\pi$, $e$, $\ln 2$, Liouville's constant, Champernowne's constant and the
  Thue--Morse sequence, in any base: the \ttx{spigot:} chance sources \\
\ttx{live} & \ttx{f1r3score}, \ttx{score-render} & off & live MIDI output
  through \ttx{midir}: the \ttx{--live PORT} option and the \ttx{ports}
  subcommand \\
\bottomrule
\end{tabularx}
\end{center}

\begin{shell}
cargo build --release --features live -p f1r3score
./target/release/f1r3score ports     # lists MIDI output ports
\end{shell}

\subsection{Build profiles}

The workspace sets \ttx{opt-level = 1} for \ttx{dev} and \ttx{opt-level = 2}
for \ttx{test}, because the conformance tests draw tens of thousands of notes.
Use \ttx{--release} for anything musical. A 16-pitch $\times$ 5-duration voice
plays at roughly 0.2--0.5\,ms per note in a release build.

\subsection{Running the tests}

\begin{shell}
cargo test --release
\end{shell}

\noindent This runs 70 tests and takes about two minutes. Nearly all of that is
\ttx{tests/law.rs} (about 60\,s) and \ttx{tests/suite.rs} (about 50\,s) in
\ttx{score-conformance}. What they cover:

\begin{center}
\small
\begin{tabularx}{\linewidth}{@{}l>{\raggedright\arraybackslash}X@{}}
\toprule
test file & what it checks \\
\midrule
\ttx{score-core} unit, \ttx{congruence.rs} & normal forms and matching; the
  congruence and round-trip property tests over wildcards (10\textsuperscript{4}
  random pairs against an independent decision procedure) \\
\ttx{score-syntax/tests/syntax.rs} & parsing, elaboration, printing and
  re-parsing, and diagnostics \\
\ttx{score-chance} unit & PRNG determinism, spigot digits against known
  expansions, interval decoding \\
\ttx{suite.rs}, \ttx{law.rs}, \ttx{t12.rs} & the ported \ttx{skeinsim} suite
  T0--T12 with the note's thresholds \\
\ttx{instruments.rs} & T13 (4-hands piano), T14 (chimera), and the instrument
  propositions \\
\ttx{outputs.rs} & records and playback, \ttx{render} against \ttx{play}, the
  MIDI golden files, replay, and detection of tampered choices \\
\ttx{examples.rs} & every file in \ttx{examples/} elaborates and plays \\
\ttx{portability.rs} & the five core crates contain no I/O \\
\bottomrule
\end{tabularx}
\end{center}

\noindent To run a single file: \ttx{cargo test --release -p score-conformance
--test instruments}.

\subsection{The cross-implementation differential}\label{ssec:differential}

The differential checks that the Rust player and the Python reference simulator
emit the same records at every resolution, null notes included. The player's
traces drive \ttx{skeinsim.py}. It needs a checkout of \ttx{publications} and
runs from the workspace root:

\begin{shell}
git clone https://github.com/F1R3FLY-io/publications.git ../publications
sh crates/score-conformance/run_differential.sh ../publications/f1r3score 40
\end{shell}

\noindent The second argument is the number of random small scores (default
40). The script ends with \ttx{differential: every scenario agrees}, or exits~1
naming each scenario that differs. It is not part of \ttx{cargo test}, because
it needs Python and the other repository.

\subsection{Building for WebAssembly}\label{ssec:wasm}

The five core crates (\ttx{score-core}, \ttx{score-logic}, \ttx{score-syntax},
\ttx{score-chance}, \ttx{score-engine}) perform no I/O. That property is
enforced by \ttx{portability.rs}, and it is what makes them candidates for
\ttx{wasm32-unknown-unknown}:

\begin{shell}
rustup target add wasm32-unknown-unknown
for c in score-core score-logic score-syntax score-chance score-engine; do
  cargo build --release --target wasm32-unknown-unknown -p $c
done
\end{shell}

\begin{caveat}{Not verified}
Neither \ttx{DESIGN.md} nor this guide has seen that build succeed, because
neither build machine had a wasm32 standard library. Ubuntu does not package
one for its \ttx{rustc}, so use rustup.
\end{caveat}

\subsection{Build problems}

\begin{description}[style=nextline,font=\normalfont\ttfamily\small]
\item[error: package `\dots` cannot be built because it requires rustc 1.75]
  The toolchain is too old; upgrade it.
\item[failed to run custom build command for `alsa-sys`]
  \ttx{--features live} on Linux without ALSA headers. Install
  \ttx{libasound2-dev} and \ttx{pkg-config}.
\item[could not compile `num-traits` \ldots{} (wasm32)] No standard library for
  that target is installed (\S\ref{ssec:wasm}).
\end{description}

%==============================================================================
\section{Deploying}\label{sec:deploy}

\subsection{Installing the binary}

\begin{shell}
cargo install --path crates/f1r3score                 # into ~/.cargo/bin
cargo install --path crates/f1r3score --features live # with live MIDI
\end{shell}

\noindent On Linux the default binary links only \ttx{libc}, \ttx{libm} and
\ttx{libgcc\_s}, so it copies between glibc systems of the same or later
vintage. It is about 15\,MB unstripped; \ttx{strip} reduces that
considerably. A fully static build (\ttx{--target x86\_64-unknown-linux-musl})
should work for the default feature set, since there are no C dependencies, but
it was not tried. The \ttx{live} build links ALSA dynamically and needs
\ttx{libasound2} at run time.

\subsection{Pin the dependency graph}

\ttx{Cargo.lock} is not committed. For a binary that promises bit-for-bit
reproducible runs it should be, so that every build resolves the same
\ttx{num-*} crates. The versions resolved for this guide were
\ttx{num-bigint}~0.4.8, \ttx{num-rational}~0.4.2, \ttx{num-traits}~0.2.19,
\ttx{num-integer}~0.1.47 and, with \ttx{live}, \ttx{midir}~0.9.1 and
\ttx{alsa}~0.7.1. Until the lock file is committed, build with the lock file
of a known-good checkout, or with \ttx{cargo build --locked} once it is.

\subsection{Reproducibility and archiving performances}

A run is determined by four things: the score, the chance sources, the
scheduler, and the matching and collection modes. Every trace records all of
them in its header, together with the score's SHA-256 structural digest and
the player version. Two properties follow for archiving:
\begin{itemize}
\item \textbf{Replay.} \ttx{f1r3score replay TRACE --score SCORE} re-runs the
  score from the trace's choices. It refuses a score whose digest differs from
  the trace's, and exits~4 if any choice emits different notes. Archive the
  score together with its trace.
\item \textbf{Playback.} \ttx{f1r3score render TRACE} needs neither the score
  nor the engine. The trace's records contain everything playback reads, so a
  trace on its own is a durable performance.
\end{itemize}
Digests are structural. Reformatting a score, renaming bound variables, or
changing stream labels (\S\ref{ssec:streams}) leaves the digest unchanged;
changing its meaning does not.

\subsection{Embedding the player as a library}

The crates can be used directly from another Rust program, for example the
F1R3Games back end or a browser front end over the core crates. This program
was compiled and run against the workspace:

\begin{lstlisting}[language=Rust]
// Cargo.toml: score-syntax, score-engine, score-chance, score-render
//             as path or git dependencies on the F1R3Score workspace
use score_engine::{Config, Engine, Limits, Sources};
use score_syntax::{load, NoFiles};

fn main() -> Result<(), String> {
    let src = std::fs::read_to_string("examples/hocket.score").map_err(|e| e.to_string())?;
    // NoFiles allows only "import std"; implement score_syntax::Loader
    // to resolve import "file.score" from wherever your scores live.
    let s = load("hocket.score", &src, &mut NoFiles).map_err(|d| d.to_string())?;
    let sources = Sources::new(score_chance::source("prng:42", 64)?);
    let mut e = Engine::new(s.arena, s.clauses, s.alph, &s.initial, s.window,
                            s.behavioural, Config::default(), sources)
        .map_err(|e| e.to_string())?;
    let limits = Limits { notes: Some(8), ..Default::default() };
    let mut notes = Vec::new();
    e.run(&limits, &mut |ev, _| notes.extend(ev.notes.iter().cloned()))
        .map_err(|e| e.to_string())?;
    let perf = score_render::performance_of(&notes, &e.alph);   // drops null notes
    print!("{}", score_render::performance_table(&perf, &e.alph));
    let midi = score_render::midi::write(&perf, &e.alph, 120);
    std::fs::write("hocket.mid", &midi.bytes).map_err(|e| e.to_string())
}
\end{lstlisting}

\noindent The pieces an embedder can use:
\begin{itemize}
\item \ttx{Sources::bind(label, source)} binds stream labels, as
  \ttx{--stream} does.
\item \ttx{Engine::step} advances one resolution at a time, for callers that
  interleave playing with something else.
\item A type implementing \ttx{score\_chance::Chance} can supply choices
  directly. A graphical chooser would be one; the CLI's \ttx{Human} type is
  the model to follow.
\end{itemize}
The engine's public types are those of revision~2 of the specification. They
are not yet a stable API: the workspace version is \ttx{0.1.0}.

\subsection{Live MIDI}

Build with \ttx{--features live}, list the available ports, and pass a
substring of the port name:

\begin{shell}
f1r3score ports
f1r3score play examples/chimera.score --live "FluidSynth" --bpm 96
\end{shell}

\noindent On a Linux host without a hardware synthesiser, a software synth
such as FluidSynth or TiMidity\texttt{++} provides an ALSA sequencer port.
Containers and headless CI machines usually lack \ttx{/dev/snd/seq}; there
\ttx{ports} fails with
\ttx{MIDI support could not be initialized} (exit~3). Prefer \ttx{--midi} there.

\begin{caveat}{Not verified}
Live MIDI compiles and links, but it has not been tested on a MIDI device.
Section~\ref{sec:limits} also describes how its timing departs from the
specification.
\end{caveat}

\subsection{Continuous integration (proposed)}

\ttx{portability.rs} and \ttx{DESIGN.md} both refer to a CI job at
\ttx{.github/workflows/ci.yml}, but that file is not in the repository. The
following workflow would cover what the design notes expect of CI: the tests,
the wasm build and the differential. It has not been run.

\begin{lstlisting}[language={}]
# .github/workflows/ci.yml  -- proposed, untested
name: ci
on: [push, pull_request]
jobs:
  test:
    runs-on: ubuntu-24.04
    steps:
      - uses: actions/checkout@v4
      - uses: dtolnay/rust-toolchain@1.75
      - run: sudo apt-get install -y libasound2-dev pkg-config
      - run: cargo build --release --features live
      - run: cargo test --release
  wasm:
    runs-on: ubuntu-24.04
    steps:
      - uses: actions/checkout@v4
      - uses: dtolnay/rust-toolchain@stable
        with: { targets: wasm32-unknown-unknown }
      - run: |
          for c in score-core score-logic score-syntax score-chance score-engine; do
            cargo build --release --target wasm32-unknown-unknown -p $c
          done
  differential:
    runs-on: ubuntu-24.04
    steps:
      - uses: actions/checkout@v4
      - uses: actions/checkout@v4
        with: { repository: F1R3FLY-io/publications, path: publications }
      - uses: dtolnay/rust-toolchain@1.75
      - run: sh crates/score-conformance/run_differential.sh publications/f1r3score 40
\end{lstlisting}

\noindent If \ttx{publications} is private, the second checkout needs a token
with read access to it.

\subsection{A container image (proposed)}

For batch rendering, such as a service that turns scores or traces into MIDI
files, a two-stage image keeps the runtime small. This Dockerfile is a
template and has not been built:

\begin{lstlisting}[language={}]
FROM rust:1.75-slim AS build
WORKDIR /src
COPY . .
RUN cargo build --release -p f1r3score

FROM debian:bookworm-slim
COPY --from=build /src/target/release/f1r3score /usr/local/bin/f1r3score
WORKDIR /work
ENTRYPOINT ["f1r3score"]
# docker run --rm -v "$PWD:/work" f1r3score play piece.score --notes 64 --midi out.mid
\end{lstlisting}

%==============================================================================
\section{Using the player}\label{sec:use}

\subsection{Quick start}

\begin{shell}
B=target/release/f1r3score
$B check  examples/twinkle.score
$B play   examples/twinkle.score --midi twinkle.mid
$B play   examples/country_voice.score --notes 32 \
          --stream A=prng:3 --stream B=prng:4 --gc freshness \
          --trace run.jsonl --midi run.mid
$B replay run.jsonl --score examples/country_voice.score
$B render run.jsonl --midi again.mid
\end{shell}

\noindent These commands, in order, check Twinkle, play it to a MIDI file,
play 32 notes of two voices while writing a trace and MIDI, replay the trace
against the score, and render the trace again without the score. The last
\ttx{again.mid} is byte-identical to \ttx{run.mid}.

\subsection{Commands}

\begin{center}
\small
\begin{tabularx}{\linewidth}{@{}>{\raggedright\arraybackslash}p{5.2cm}>{\raggedright\arraybackslash}X@{}}
\toprule
command & purpose \\
\midrule
\ttx{check SCORE} & Parse, elaborate and lint. Prints the digest, the regime
  (\ttx{local}, or \ttx{window-behavioural (depth $d$)}; see
  \S\ref{ssec:fragment}), and the number of terms and clauses. \\
\ttx{expand SCORE} & Print the elaborated core term, with every definition
  expanded. The output re-parses. \\
\ttx{play SCORE [opts]} & Run the score once (\S\ref{ssec:playopts}). \\
\ttx{replay TRACE --score SCORE} & Re-run from a trace's recorded choices and
  check that each emits the recorded notes. Accepts \ttx{--midi}, \ttx{--bpm}
  and \ttx{--format}; the scheduler, mode and collection come from the
  trace. \\
\ttx{render TRACE [opts]} & Playback only: read the notes off the trace's
  records. Accepts \ttx{--midi}, \ttx{--bpm}, \ttx{--out}, \ttx{--format} and
  \ttx{--include-null}. \\
\ttx{machine import FILE --kind pitch|dur [--dual] --score SCORE [--name N]} &
  Convert a machine exported by the instrument (indexed states) into a
  \ttx{factor} line over the score's alphabet. \\
\ttx{ports} & \ttx{live} builds only: list MIDI output ports. \\
\bottomrule
\end{tabularx}
\end{center}

\noindent \ttx{check} on Twinkle prints:
\begin{out}
ok examples/twinkle.score
digest  d5b61ae1d112ed6322f0981a125f2d82ce54cfea3277e89636bda081845af0a9
regime  local
terms   163   clauses 1
\end{out}

\subsection{Options for \ttx{play}}\label{ssec:playopts}

\begin{center}
\small
\begin{longtable}{@{}l>{\raggedright\arraybackslash}p{9.6cm}@{}}
\toprule
option & meaning (default) \\
\midrule
\endhead
\ttx{--chance SPEC} & default chance source for every stream not bound
  otherwise (\ttx{prng:0}) \\
\ttx{--stream LABEL=SPEC} & bind a stream label, a timbre, or
  \ttx{LABEL.pitch}\,/\,\ttx{LABEL.dur} to a source; repeatable \\
\ttx{--scheduler S} & order of simultaneously due sets: \ttx{canonical},
  \ttx{random:SEED}, or \ttx{by-timbre:A,B,...} (\ttx{canonical}) \\
\ttx{--mode max|one} & fire a maximal matching, or a single candidate as a
  diagnostic (\ttx{max}) \\
\ttx{--gc off|freshness} & collect inert record locations (\ttx{off}) \\
\ttx{--notes N} & stop after N notes of positive length \\
\ttx{--until T} & stop at musical time T in whole notes, e.g.\ \ttx{3/2} \\
\ttx{--max-steps N} & stop after N resolutions \\
\ttx{--max-alternatives N} & bound on maximal matchings per set (100000) \\
\ttx{--unproductive N} & bound on consecutive null-note steps (100000) \\
\ttx{--max-states N} & bound on behavioural window graphs (200000) \\
\ttx{--kmax K} & interval-decoding digit bound per choice (64) \\
\ttx{--trace FILE} & write a replayable JSONL trace \\
\ttx{--format table|json} & performance on stdout (\ttx{table}) \\
\ttx{--out FILE} & performance as JSON to FILE \\
\ttx{--include-null} & include null notes in the printed performance \\
\ttx{--midi FILE}, \ttx{--bpm N} & Standard MIDI File, at N quarter notes per
  minute (120) \\
\ttx{--live PORT} & live MIDI to the first port whose name contains PORT
  (\ttx{live} builds) \\
\ttx{--quiet} & print no performance \\
\bottomrule
\end{longtable}
\end{center}

\noindent With none of \ttx{--notes}, \ttx{--until} or \ttx{--max-steps}, the
player warns that an unbounded score will play forever. Written scores such as
Twinkle stop by themselves when nothing is left to communicate. Voices do not:
always bound them.

\subsection{Chance sources}\label{ssec:chance}

Every automatic source is a stream of digits, and every choice is made by exact
interval decoding. No floating point enters a weight or a choice, so the same
sources yield the same run on every platform.

\begin{center}
\small
\begin{tabularx}{\linewidth}{@{}l>{\raggedright\arraybackslash}X@{}}
\toprule
spec & source \\
\midrule
\ttx{prng:SEED} & SplitMix64 feeding xoshiro256**, fixed in the source so that
  it cannot drift with a dependency \\
\ttx{spigot:C/B} & digits of the fractional part of constant C in base B, with
  C one of \ttx{pi}, \ttx{e}, \ttx{ln2}, \ttx{liouville}, \ttx{champernowne},
  \ttx{thue-morse} \\
\ttx{spigot:C/B@Q} & the same, starting at digit Q. Use this when the first
  note is written rather than drawn. \\
\ttx{human} & an interactive chooser on the terminal (\S\ref{ssec:human}) \\
\bottomrule
\end{tabularx}
\end{center}

\noindent A set with a single alternative, which includes every null note,
consumes no digits.

\subsection{Stream labels}\label{ssec:streams}

Choices are routed to sources through \emph{stream labels}. Independent voices
should draw from independent streams: the note's product-law theorem holds
surely only then.
\begin{itemize}
\item \ttx{Voice(...) \#A} labels every receipt in that application's expansion
  \ttx{A}, and continuations inherit the label of the receipt that fired them.
  Labels are metadata: they do not change the digest.
\item An unlabelled receipt takes its subject's timbre as its label, or
  \ttx{loc:<digest>} when the timbre is open. So \ttx{--stream vibes=human}
  addresses every unlabelled receipt in timbre \ttx{vibes}.
\item A contention set draws from the stream of its least receipt occurrence.
  At a chimera that is the first-installed timbre's stream, so a chimera is
  normally given one explicit label.
\item \emph{Factor routing.} \ttx{--stream A.pitch=... --stream A.dur=...}
  resolves a voice's pitch and duration by two independent draws (Lemma~6.4 of
  the note). This happens only when both are bound, the set is a star, and each
  top-level factor reads only one side. Otherwise the draw is joint. The trace
  records which happened (\ttx{"routing":"factored"} or \ttx{"joint"}).
\end{itemize}

\subsection{Schedulers, modes and collection}

\paragraph{Schedulers.} When several contention sets are due at the same
onset, the scheduler orders them: \ttx{canonical} is a fixed order,
\ttx{random:SEED} a seeded shuffle, and \ttx{by-timbre:A,B} resolves timbre A's
sets first. Independent voices give the same performance under every
scheduler. For example, Twinkle's \ttx{--out} files under
\ttx{by-timbre:bass,piano} and \ttx{random:9} are byte-identical.

\paragraph{Modes.} \ttx{max} is the semantics. \ttx{one} fires a single
candidate per set and exists only for diagnosis (test T3).

\paragraph{Collection.} Voices leave behind record locations whose alternatives
went unplayed. They are inert, but they sit in the soup. \ttx{--gc freshness}
removes record-shaped locations that no longer have an available candidate. It
never changes a performance, and it guards its own soundness: spawning at a
collected location is the run-time error \ttx{freshness-violated}. Use it for
long runs of voices. It collects nothing on instruments.

\subsection{Human choice}\label{ssec:human}

\ttx{--chance human}, or \ttx{--stream LABEL=human} for one stream, stops at
every set with more than one alternative. It lists the alternatives as timbre,
pitch, duration and carried next datum (after \ttx{->}), with their normalised
probabilities for guidance, and reads an index from standard input:

\begin{out}
step 1 at onset 0 (stream marimba):
  [0] marimba E4 q -> C4   (0.100)
  [1] marimba E4 q -> D4   (0.100)
  ...
  [9] marimba E4 e -> G4   (0.100)
choose 0..9>
\end{out}

\noindent Only alternatives of positive weight are offered, so every trace a
human produces is admissible. The chooser is step-wise, not real time. Binding
\ttx{human} to one label and PRNGs to the rest gives a mixed performance, with
a person improvising against algorithmic players. If input ends while a choice
is pending, the run stops with exit~3, after printing what was played.

\subsection{Outputs}

\paragraph{Performance.} This is the multiset of timed notes of positive
length, sorted by onset, timbre, pitch and duration, so that equal performances
produce byte-identical files. Times and lengths are exact rationals in whole
notes. The table form:

\begin{out}
     onset  timbre     pitch  dur     length
         0  piano      C4     q          1/4
         0  bass       C3     h          1/2
       1/4  piano      C4     q          1/4
\end{out}

\noindent The JSON form (\ttx{--format json}, or \ttx{--out FILE}):
\begin{out}
{"score":"sha256:aa81...d0e8","end":"5/4","notes":[
  {"onset":"0","timbre":"piano","pitch":"A2","dur":"h","len":"1/2"}, ... ]}
\end{out}

\paragraph{Trace.} JSON Lines, with format \ttx{f1r3score-trace/2}. The first
line is a header:
\begin{itemize}
\item the player version, the score path and digest;
\item the alphabets, with MIDI numbers and lengths, and the timbres;
\item every source and stream binding, the scheduler, mode, locality regime,
  collection mode, and \ttx{kmax}.
\end{itemize}
Then comes one line per resolution:
\begin{out}
{"step":1,"onset":"0","loc":"31933326af6fe475","set":2,"alts":2,"chosen":0,
 "routing":"joint","digits":{"B":1},
 "records":[{"onset":"0","loc":"31933326af6fe475","recv":["bass","C4"],
             "send":["bass","h"],"payload":["8307dfa7ddbd16d9","bass","E4"]}]}
\end{out}
\noindent Each record gives the receipt's and the message's subject
decorations and the payload decoration, with the payload process given by its
digest. Notes are not stored; they are computed from records.

\paragraph{MIDI.} \ttx{--midi} writes a format-1 Standard MIDI File, with a
tempo track and one track per rendered timbre. Each track carries a program
change from the timbre's \ttx{gm(N)} and the channel from its \ttx{on C}. A
note's track is chosen by its own timbre, so a chimera's notes are routed note
by note to whichever timbre won. Rests and notes in the reserved timbres
produce no events. The division is the least number of ticks per quarter that
makes every onset and length integral, if that is at most 32767. Otherwise it
is 960, with rounding and a warning. \ttx{--bpm} is a rendering parameter only.

\subsection{Exit codes and stop reasons}

\begin{center}
\small
\begin{tabular}{@{}cl@{}}
\toprule
code & meaning \\
\midrule
0 & success \\
1 & usage error (unknown option, missing value, missing file) \\
2 & parse or elaboration error (and malformed traces for \ttx{render}) \\
3 & run-time error: \ttx{unproductive}, \ttx{freshness-violated}, alternative
  limit, state bound, I/O \\
4 & replay mismatch: digest differs, a choice emits different notes, or
  choices are left over \\
\bottomrule
\end{tabular}
\end{center}

\noindent On success the player reports on standard error why it stopped:
\ttx{Quiescent} (nothing left to communicate), \ttx{Notes}, \ttx{Until} or
\ttx{Steps}. Examples are \ttx{Quiescent after 22 steps, 22 notes} and
\ttx{Notes after 96 steps, 32 notes}. The performance goes to standard output
and diagnostics to standard error, so \ttx{--format json > perf.json} is safe.

%==============================================================================
\section{Writing scores}\label{sec:write}

\subsection{File structure}

A score file is UTF-8, with comments \ttx{//} and \ttx{/* */}. It contains, in
any order, an optional \ttx{score NAME}, imports, the three alphabet
declarations, \ttx{pitchset}, \ttx{def} and \ttx{factor} items, and exactly one
\ttx{play} item giving the initial configuration.

\begin{lstlisting}[language=score]
score First
import std
pitches   { r, C4, E4, G4, C5 }
durations { q = 1/4, h = 1/2 }
timbres   { piano = gm(0) on 1 }

// a melody over a sustained two-note chord
play line(base "Mel", piano) [ C4 q, E4 q, G4 q, C5 q ]
   | line(base "Lo", piano)  [ C4 h, C4 h ]
   | line(base "Hi", piano)  [ G4 h, E4 h ]
\end{lstlisting}

\noindent Parallel lines sound together; that is harmony as independence at
composition. The run:
\begin{out}
     onset  timbre     pitch  dur     length
         0  piano      C4     q          1/4
         0  piano      C4     h          1/2
         0  piano      G4     h          1/2
       1/4  piano      E4     q          1/4
       1/2  piano      C4     h          1/2
       1/2  piano      E4     h          1/2
       1/2  piano      G4     q          1/4
       3/4  piano      C5     q          1/4
f1r3score: Quiescent after 8 steps, 8 notes
\end{out}

\noindent \ttx{import std} brings in the standard library. \ttx{import
"lib.score"} reads a file relative to the importing score; an imported file may
contain only \ttx{def}, \ttx{factor}, \ttx{pitchset} and \ttx{import}.

\subsection{Alphabets}

\paragraph{Pitches.} \ttx{pitches \{ ... \}} lists the pitch alphabet \emph{in
order}; that order defines \ttx{step(k)}. The rest \ttx{r} is required and is
unordered. A pitch in scientific notation (\ttx{C4}, \ttx{F\#3}, \ttx{Bb2}) gets
its MIDI number automatically, with C4 = 60. Any other identifier needs one:
\ttx{deg7 = 67}. Enharmonics are distinct data. Two abbreviations are
available:
\begin{itemize}
\item \ttx{scale KIND LO .. HI}, inclusive, where KIND is one of \ttx{major}
  (\ttx{ionian}), \ttx{minor} (\ttx{aeolian}), \ttx{dorian}, \ttx{pentatonic}
  (\ttx{major-pentatonic}), \ttx{minor-pentatonic}, \ttx{blues} or
  \ttx{chromatic}. A flat low pitch (\ttx{Bb3}) names the scale with flats;
\item \ttx{piano88}, which gives A0 to C8.
\end{itemize}
\ttx{pitchset NAME = \{...\}} declares a named subset for use in definitions.

\paragraph{Durations.} Each duration is \ttx{NAME = RATIONAL} in whole notes:
\ttx{q = 1/4}, \ttx{t = 1/6} for a triplet eighth, \ttx{dq = 3/8} for a dotted
quarter. The null duration \ttx{eps = 0} is implicit and may not be
redeclared. Pitch and duration names must be disjoint.

\paragraph{Timbres.} \ttx{NAME = gm(PROGRAM) on CHANNEL}, with a General MIDI
program from 0 to 127 and a channel from 1 to 16. The channel defaults to the
timbre's position, modulo 16. \ttx{dead}, \ttx{ctl} and \ttx{\_} are reserved.

\subsection{Processes}

\begin{center}
\small
\begin{tabularx}{\linewidth}{@{}l>{\raggedright\arraybackslash}X@{}}
\toprule
surface & meaning \\
\midrule
\ttx{0} & the inert process \\
\ttx{<@P, tau, d>} & a name; \ttx{\_} for an open timbre or datum.
  Compound quotes are parenthesised: \ttx{<@(P | Q), piano, q>} \\
\ttx{@P} & the general name \ttx{<@P, \_, \_>} \\
\ttx{x!(Q, tau, d)} & a message whose payload is the name \ttx{<@Q, tau, d>} \\
\ttx{x!(Q)} & a message passing the general name \ttx{@Q} \\
\ttx{for (y <- x where psi) \{ P \}} & a receipt; the clause is optional
  (weight 1) \\
\ttx{for (y1 <- x1 \& y2 <- x2 where psi) \{ P \}} & a join (chord receipt) \\
\ttx{P | Q} & parallel composition \\
\ttx{*x} & drop: run the process quoted by \ttx{x} \\
\ttx{at 3/4 \{ P \}} & an entrance: stamp P at 3/4 (top level of \ttx{play}
  only; the default is 0) \\
\ttx{base "tag"} & a closed process used as a location; distinct tags give
  inequivalent bases \\
\ttx{keyloc(P, k)}, \ttx{keycode(P, k, tau)} & a key's location and its code
  location on keyboard base P \\
\bottomrule
\end{tabularx}
\end{center}

\noindent In binder position \ttx{\_} is an unused binder. In decoration
position it is the wildcard.

\subsection{Definitions}

Definitions are macros, expanded at elaboration time:
\begin{lstlisting}[language=score]
def N(L: proc, tau: timbre, p: pitch, d: dur; K: proc) =
  for (_ <- <@L, tau, p>) { K } | <@L, tau, d>!(0)
def P1 = base "P1"                                        // no parameters
def touch(P: proc, k: pitch, d: dur): name = <@keyloc(P, k), _, d>   // returns a name
\end{lstlisting}

\begin{itemize}
\item \textbf{Sorts.} The parameter sorts are \ttx{proc}, \ttx{name},
  \ttx{timbre}, \ttx{pitch}, \ttx{dur}, \ttx{clause}, \ttx{string},
  \ttx{pitchset}, lists \ttx{[S]}, tuples \ttx{(S1, S2)}, and \ttx{cont}.
  A \ttx{cont} is a definition of arity one in a name, passed by name, with
  \ttx{\_} marking the hole: \ttx{Req(\_, tau, BS)}.
\item \textbf{No recursion.} The call graph must be acyclic. A recursive
  definition is rejected with a pointer to \S5.2 of the note: unbounded
  behaviour comes from reflection, as in \ttx{Voice}.
\item \textbf{Comprehensions.} \ttx{par v in durations+, t in pitches \{ ...
  \}} expands to a parallel composition. The generators are \ttx{pitches},
  \ttx{pitches-r}, \ttx{durations}, \ttx{durations+} (positive durations),
  declared pitch sets, literal sets, and lists, which may be destructured as
  \ttx{par (tau, psi) in voices}.
\item \textbf{Lines.} \ttx{line(L, tau) [ C4 q, E4 h, ... ]} is a
  right-nested chain of written notes. \ttx{then K} replaces the final
  \ttx{0}, for handing a phrase to a responder.
\item \textbf{Synchronous output.} \ttx{a ; b ; 0} runs \ttx{b} when
  \ttx{a}'s note has ended. A step is one touch or a parenthesised parallel
  group. \ttx{;} binds more tightly than \ttx{|}, so \ttx{a ; 0 | c ; 0} is two
  players. A \ttx{;} inside a message payload is an error. Inside definition
  arguments \ttx{;} separates arguments, so parenthesise a sequence there.
\item \textbf{Strings.} \ttx{++} concatenates strings, as in
  \ttx{base tag ++ "/S"}.
\end{itemize}

\subsection{Clauses}\label{ssec:clauses}

A clause is a \emph{graded} formula with values in the non-negative rationals.
It is built from factors, which can be named with \ttx{factor NAME = ...}.

\begin{center}
\small
\begin{tabularx}{\linewidth}{@{}l>{\raggedright\arraybackslash}X@{}}
\toprule
graded form & value \\
\midrule
\ttx{3/2}, \ttx{0.3}, \ttx{20} & a constant; decimals are exact (0.3 is 3/10) \\
\ttx{[b]} & 1 if the crisp formula b holds, else 0 \\
\ttx{psi * psi}, \ttx{psi + psi} & product (short-circuits on 0), sum \\
\ttx{table KEY \{ k: w, ..., \_: w \}} & lookup by the listed key; \ttx{\_} is the
  default and an absent default is 0. Keys are \ttx{dur}, \ttx{step} (ranges
  such as \ttx{-1..1} allowed, but they must not overlap), \ttx{(pitch,
  carry)}, \ttx{(prev, dur)}, or tuples of view functions \\
\ttx{machine pitch \{ C4 -> E4 @ 2, ... \}} & a table keyed by (pitch, carry):
  a weighted automaton on pitches \\
\ttx{machine dur \{ q -> e @ 1, ... \}} & a table keyed by (prev, dur) \\
\ttx{machine dur uniform}, \ttx{machine pitch uniform} & the complete uniform
  machine \\
\ttx{machine pitch dual \{...\}} & the dual (rhythm-leads) reading \\
\ttx{machine pitch import "f.machine"} & a machine exported by the instrument \\
\ttx{sum x in S \{ psi \}} & a sum expanded at elaboration \\
\bottomrule
\end{tabularx}
\end{center}

\noindent Crisp formulae have \ttx{true}, \ttx{false}, \ttx{!}, \ttx{\&\&},
\ttx{||}, the comprehensions \ttx{any x in S \{ b \}} and \ttx{all x in S \{ b
\}}, and these atoms:

\begin{center}
\small
\begin{tabularx}{\linewidth}{@{}l>{\raggedright\arraybackslash}X@{}}
\toprule
atom & holds when \\
\midrule
\ttx{pitch(p)}, \ttx{dur(v)} & the candidate's note has this pitch, this
  duration \\
\ttx{carry(e)} & the datum carried to the next step is e; \ttx{carry(\_)} when
  it is open \\
\ttx{step(k)} & carried and current pitch are ordered and differ by k
  positions in the pitch alphabet \\
\ttx{held(t)}, \ttx{prev(u)} & the location's last record carries pitch t;
  the previous duration was u \\
\ttx{back(j, e)} & the record j levels into the past carries e \\
\ttx{last[m1 .. mk]}, \ttx{call[...]} & the last k notes match; the same on
  a handed-over call \\
\ttx{at(s)}, \ttx{passes(s)} & the location's quote, or the passed process,
  satisfies the spatial pattern s \\
\ttx{pitch\#i}, \ttx{dur\#i}, \ttx{carry\#i} & per-pattern atoms of a chord
  receipt \\
\ttx{leads(b)} & the configuration after firing the candidate satisfies b
  (\S\ref{ssec:fragment}) \\
\bottomrule
\end{tabularx}
\end{center}

\noindent Spatial patterns are \ttx{<n, tau, d>!(s, tau, d)} for a message,
\ttx{for(<n, tau, d>) s} for a receipt, \ttx{0}, \ttx{true}, \ttx{s | s},
\ttx{!s}, \ttx{s \&\& s} and \ttx{@s}. In patterns, \ttx{?} matches any
component, \emph{including} a wildcard, while \ttx{\_} matches only a
wildcard. A location is never a wildcard, so write \ttx{<?, ...>} rather than
\ttx{<\_, ...>}.

\paragraph{Encourage or enforce.} Crisp factors enforce an idiom, and graded
factors encourage one. Two examples from the conformance scores:
\begin{lstlisting}[language=score]
factor rule  = [!(any t in pitches { back(0, t) && back(1, t) && carry(t) })]  // no triple repeat
factor boost = [back(1, C4) && back(0, D4) && carry(E4)] * 20 + 1             // favour C-D-E
\end{lstlisting}

\subsection{The checkable fragment}\label{ssec:fragment}

A clause is admitted if it is in one of two classes.
\begin{itemize}
\item \textbf{Local.} No behavioural formula appears. The clause's value
  depends only on the candidate.
\item \textbf{Window-behavioural.} \ttx{leads(phi)} appears, where
  \ttx{phi} is built from \ttx{true}, \ttx{false}, \ttx{live}, \ttx{!},
  \ttx{\&\&}, \ttx{||}, \ttx{<g> phi}, \ttx{<> phi}, \ttx{nu X. phi} and
  \ttx{X}. Every clause in the score must read the past to a bounded depth. The
  player then decides \ttx{phi} exactly on a finite graph of configurations,
  truncated to that depth.
\end{itemize}
\ttx{check} reports which regime a score uses. Graded behavioural modalities
and graded fixed points are rejected with a reason. A bare modality must go
under \ttx{leads(...)}. Exceeding \ttx{--max-states} is an error, never a
truncation. The use for all this is viability, meaning enforcement that never
paints a voice into a corner:
\begin{lstlisting}[language=score]
factor viable = [leads(nu X. <!(any t in pitches { back(0, t) && back(1, t) && carry(t) })> X)]
play Voice(piano, style * rhythm * rule * viable, C4, q, "A") #A
\end{lstlisting}

\subsection{The standard library}

\begin{center}
\small
\begin{longtable}{@{}>{\raggedright\arraybackslash}p{5.3cm}>{\raggedright\arraybackslash}p{9.4cm}@{}}
\toprule
definition & what it is \\
\midrule
\endhead
\ttx{N(L, tau, p, d; K)}, \ttx{line} & a written note, and a written line \\
\ttx{Rec}, \ttx{Fan(x, tau, over)}, \ttx{Seed} & the record, the fan of all
  continuations (over a pitch set), and the virtual first record \\
\ttx{Step}, \ttx{StepOver} & one step of a voice: every key offered, the
  clause keeps the live one \\
\ttx{Req}, \ttx{Refresh}, \ttx{Server}, \ttx{ServerTo} & the reflective
  machinery: request, code refresh, server \\
\ttx{Voice(tau, psi, p0, u0, tag)} & an unbounded voice with clause
  \ttx{psi}, first pitch \ttx{p0}, previous duration \ttx{u0}. \ttx{VoiceOver}
  takes a pitch set \\
\ttx{DualVoice} & rhythm leads: the duration is chosen a note early and
  carried \\
\ttx{Hocket(tauA, psiA, tauB, psiB, p0, u0, tag)} & two players sharing one
  line, note for note, possibly on different instruments \\
\ttx{Call}, \ttx{Resp}, \ttx{Responder}, \ttx{RespAny}, \ttx{ResponderAny} &
  call and response. The call hands over its whole history; \ttx{...Any}
  answers in whatever timbre the caller chose \\
\ttx{Key}, \ttx{Installed}, \ttx{KeyRefresh} & self-re-arming keys \\
\ttx{Keyboard(P, tau, psi, over)} & every key of a pitch set on base \ttx{P} \\
\ttx{Chimera(P, k, voices)}, \ttx{ChimeraKeyboard(P, over, voices)} & one key
  location holding a key per \ttx{(timbre, clause)} \\
\ttx{touch(P, k, d)} & a touch name with the timbre left to the instrument;
  write \ttx{touch(P, A3, q)!(0)} \\
\bottomrule
\end{longtable}
\end{center}

\noindent Read \ttx{crates/score-syntax/std.score} for the definitions
themselves. It is short and fully commented. \ttx{f1r3score expand} shows what
any use of it elaborates to.

%==============================================================================
\section{Worked examples}\label{sec:examples}

\subsection{A voice with a style}

A voice draws each note from its clause. Here the clause is mostly stepwise
motion, a small rhythm machine, and no rests:
\begin{lstlisting}[language=score]
score Walk
import std
pitches   { r, scale major C4 .. C5 }
durations { e = 1/8, q = 1/4, h = 1/2 }
timbres   { piano = gm(0) on 1 }

factor steps  = table step { -1..1: 4, -2: 1, 2: 1, _: 0 }   // mostly stepwise
factor rhythm = machine dur { q -> q @ 3, q -> e @ 1, e -> e @ 1, e -> q @ 2,
                              q -> h @ 1, h -> q @ 1 }
factor norest = [!carry(r)]

play Voice(piano, steps * rhythm * norest, C4, q, "W") #W
\end{lstlisting}
\begin{shell}
f1r3score play walk.score --notes 12 --stream W=prng:5 --gc freshness --midi walk.mid
\end{shell}
\begin{out}
     onset  timbre     pitch  dur     length
         0  piano      C4     q          1/4
       1/4  piano      C4     q          1/4
       1/2  piano      D4     q          1/4
       3/4  piano      E4     h          1/2
       5/4  piano      D4     q          1/4
       ...
f1r3score: Notes after 36 steps, 12 notes
\end{out}
\noindent Each note takes three resolutions: one note and two null notes of
the reflective server. Change the seed to get another performance of the same
score. Change one factor to move the style to a new instrument.

\subsection{The shipped examples}

\begin{center}
\small
\begin{tabularx}{\linewidth}{@{}l>{\raggedright\arraybackslash}X>{\raggedright\arraybackslash}p{4.4cm}@{}}
\toprule
file & what it shows & try \\
\midrule
\ttx{twinkle.score} & written lines and harmony by independence; golden score
  of T12 & \ttx{play ... --midi t.mid} \\
\ttx{country\_voice.score} & a product of experts, style $\times$
  physicality $\times$ rhythm, against a walking bass & \ttx{--notes 32
  --stream A=prng:3 --stream B=prng:4 --gc freshness} \\
\ttx{hocket.score} & two players sharing a line, marimba and vibes
  alternating & \ttx{--notes 16} \\
\ttx{call\_response.score} & a flute calls and a cello answers a fifth up
  after reading the call & \ttx{--notes 8} \\
\ttx{four\_hands.score} & two keyboards and two players with synchronous
  output; golden score of T13 & \ttx{--include-null} shows the 16 null notes \\
\ttx{chimera.score} & detached eighths sound on vibes and sustained halves on
  strings; the instrument chooses its timbre & \ttx{--midi c.mid} \\
\bottomrule
\end{tabularx}
\end{center}

\noindent Call and response, for instance, plays:
\begin{out}
         0  flute      C4     q          1/4
       1/4  flute      E4     q          1/4
       1/2  flute      G4     q          1/4
       3/4  cello      G4     h          1/2
       5/4  cello      D5     h          1/2
       7/4  cello      C5     h          1/2
\end{out}

\subsection{Transcendental numbers as players}

The instrument's embedding of $\pi$ and $e$ is itself a score. Route the
voice's pitch factor to the hexadecimal digits of $\pi$ and its duration factor
to the base-5 digits of $e$:
\begin{shell}
f1r3score play crates/score-conformance/scores/t4_embedding.score --notes 8 \
  --gc freshness --stream A.pitch=spigot:pi/16@1 --stream A.dur=spigot:e/5
\end{shell}
\noindent The pitches after the written D4 are E4, D\#4, r, F\#4, A\#4, G\#4,
G\#4. On the 16-state alphabet C4 \dots{} D5, r, those are the hex digits 4,
3, F, 6, A, 8, 8 of $\pi$. \ttx{@1} skips the first digit, because the first
pitch is written, not drawn.

\subsection{Importing a drawn machine}

A weighted automaton drawn in the F1R3Skein instrument is exported with indexed
states:
\begin{lstlisting}[language={}]
machine 3 start 1 { 0 -> 1 @ 1, 1 -> 1 @ 3, 1 -> 2 @ 2, 2 -> 1 @ 2, 2 -> 2 @ 1 }
\end{lstlisting}
\begin{shell}
f1r3score machine import dur.machine --kind dur --score examples/country_voice.score --name groove
\end{shell}
\begin{out}
factor groove = machine dur { h -> q @ 1, q -> q @ 3, q -> e @ 2, e -> q @ 2, e -> e @ 1 }
\end{out}
\noindent The states map onto the score's alphabet by position. A duration
machine uses the positive durations in declaration order. A pitch machine uses
the ordered pitches followed by the rest. The state count must equal the
alphabet size.

%==============================================================================
\section{Troubleshooting}\label{sec:trouble}

\begin{description}[style=nextline,font=\normalfont\ttfamily\small]
\item[error: `D4` is not a declared pitch or duration]
  Every datum and timbre must be declared. Exit 2.
\item[error: recursive definition: `A` refers to `A` \dots]
  Use reflection (a \ttx{Server} with its code on a code channel), as
  \ttx{Voice} does.
\item[error: `dead` is the reserved dead timbre \dots]
  Only \ttx{base}, \ttx{keyloc} and \ttx{keycode} may produce it.
\item[error: a behavioural modality is a formula of the configuration: write it under `leads(...)`]
  A bare \ttx{<>} in a clause.
\item[error: a table key is listed twice]
  Overlapping ranges in a \ttx{step} table, such as \ttx{-1..1} and
  \ttx{-2..2}. List the outer values separately.
\item[error: expected a process, found `1`]
  Only \ttx{0} is a process literal; use \ttx{base "..."} for a fresh
  location.
\item[warning: a subject has datum `\_` \dots{} can never communicate]
  A name with no polarity. The score still runs, but that receipt or message
  is dead.
\item[warning: a top-level receipt and message share a location but their subjects never match]
  Same polarity, or timbres that cannot meet (both open, or different).
\item[Quiescent after 8 steps, 2 notes (with a larger --notes)]
  The score deadlocked: every candidate has weight 0. A crisp rule has
  painted a voice into a corner. Add a viability factor
  (\S\ref{ssec:fragment}), or soften the rule into a graded factor. Compare
  \ttx{viable\_absent.score} with \ttx{viable.score} in the conformance
  scores.
\item[error: unproductive]
  Too many consecutive null-note steps: a reflective loop that never plays.
  Exit 3.
\item[the run never ends]
  A voice is unbounded. Give \ttx{--notes}, \ttx{--until} or
  \ttx{--max-steps}.
\item[memory grows over a very long run]
  Terms are hash-consed and never freed; each note interns about 80 fan
  records. \ttx{--gc freshness} trims the soup but not the arena. Split very
  long performances.
\item[replay-mismatch: step N: the replayed choice emitted different notes]
  The trace was edited, or the score changed in a way that keeps the digest
  but not the choices. That should not happen, so report it with both files.
\item[digest mismatch: the trace is of score \dots]
  Replay against the score the trace was made from.
\end{description}

%==============================================================================
\section{Known limitations and discrepancies}\label{sec:limits}

The following were found while preparing this guide, or are recorded in
\ttx{DESIGN.md}. None affects the correctness of offline runs.

\begin{enumerate}
\item \textbf{No CI workflow in the repository.} \ttx{portability.rs} and
  \ttx{DESIGN.md} refer to \ttx{.github/workflows/ci.yml}, which does not
  exist. As a result, the wasm32 build and the differential are not run
  automatically. A proposed workflow is in \S\ref{sec:deploy}.
\item \textbf{wasm32 is unverified.} The core crates are I/O-free by test,
  but their wasm build has not been seen to succeed.
\item \textbf{\ttx{Cargo.lock} is not committed}, which weakens the
  reproducibility the design otherwise guarantees.
\item \textbf{Live MIDI timing differs from the specification.} The
  specification has a playback thread fed through a bounded channel, with the
  engine running ahead by a look-ahead and never waiting on the clock. The
  implementation instead sleeps inside the engine's callback until each
  note-on is due, so the engine waits on the clock. Live MIDI has not been
  tested on a device.
\item \textbf{\ttx{--live} requires PORT}, matched as a substring, whereas the
  specification makes it optional. The \ttx{ports} subcommand exists only in
  \ttx{live} builds and is not listed in the usage text.
\item \textbf{Lints print as \ttx{error: warning: \dots}}, even though
  elaboration succeeds and the exit code is 0.
\item \textbf{A stale \ttx{replay:FILE} hint.} The error for an unknown chance
  source suggests \ttx{replay:FILE}, which is not implemented as a source.
  Replay is the \ttx{replay} subcommand.
\item \textbf{Single-alternative events are not checked on replay.} Their
  \ttx{chosen} field is not validated, since nothing was chosen. Tampering
  with a real choice is detected.
\item \textbf{Command-line forms differ from the specification.} The
  specification writes \ttx{replay piece.score run.jsonl} and
  \ttx{machine import FILE --pitches SCORE}. The binary takes \ttx{replay TRACE
  --score SCORE} and \ttx{machine import FILE --kind pitch|dur --score SCORE}.
  This guide follows the binary.
\item \textbf{Deadlock looks like a normal end.} It ends as \ttx{Quiescent}
  with exit 0. Compare the note count with \ttx{--notes} to detect it.
\item \textbf{Scope.} The following are outside this version:
  \begin{itemize}
  \item graded behavioural modalities and graded fixed points;
  \item atoms reading the note's timbre (reserved for orchestration);
  \item a real-time human chooser;
  \item tempo as a graded dimension.
  \end{itemize}
\item \textbf{\ttx{spigot\_stream} is not used.} The upstream crate in
  F1R3Games had defects outside base 10 (\ttx{DESIGN.md}), so
  \ttx{score-chance} computes its own digits.
\end{enumerate}

%==============================================================================
\appendix
\section{Repository layout}

\begin{center}
\small
\begin{tabularx}{\linewidth}{@{}l>{\raggedright\arraybackslash}X@{}}
\toprule
path & contents \\
\midrule
\ttx{crates/score-core} & exact rationals, alphabets, decorations with
  wildcards, hash-consed terms in congruence normal form, substitution,
  matching, records and $\nu$, bases and key locations, SHA-256 digests \\
\ttx{crates/score-logic} & spatial, crisp and graded formulae, tables, the
  checkable fragment, candidate views, evaluation \\
\ttx{crates/score-syntax} & lexer, parser, elaborator, printer, \ttx{std.score} \\
\ttx{crates/score-chance} & PRNG, spigot digits, exact interval decoding \\
\ttx{crates/score-engine} & soup, candidates, contention sets, maximal
  matchings, time, schedulers, collection, factor routing, the behavioural
  explorer \\
\ttx{crates/score-render} & playback, performance and trace JSON, trace
  parsing, MIDI writer, live MIDI \\
\ttx{crates/f1r3score} & the command-line player \\
\ttx{crates/score-conformance} & tests T0--T14, the scores and golden MIDI
  files they use, \ttx{gen\_scores.py}, \ttx{differential.py},
  \ttx{run\_differential.sh} \\
\ttx{examples/} & six example scores \\
\ttx{DESIGN.md} & decisions, deviations, conformance results, what is and is
  not verified \\
\ttx{docs/} & \emph{Scores as Processes} (PDF) \\
\bottomrule
\end{tabularx}
\end{center}

\noindent The five core crates depend only on the \ttx{num-*} crates and
perform no I/O. The licence is Apache~2.0.

\section{Conformance tests}

\begin{center}
\small
\begin{tabularx}{\linewidth}{@{}l>{\raggedright\arraybackslash}X@{}}
\toprule
test & claim checked \\
\midrule
T0 & inertness of unplayed alternatives \\
T1 & a voice's law is the normalised product of its factors (20{,}000 notes) \\
T2 & independent voices: interleaving invariance, product law \\
T3 & a note cannot be fitted from both sides; \ttx{max} vs \ttx{one} mode \\
T4 & the $\pi$/$e$ embedding, by factor routing \\
T5 & the dual (rhythm-leads) voice \\
T6 & simultaneity by contention \\
T7 & style $\times$ physicality \\
T8 & motifs: forbid (crisp), encourage (graded) \\
T9 & tables equal their formulae \\
T10, T11 & reflective voice equals unrolled voice; inertness under reflection \\
T12 & Twinkle: 22 notes ending at 4, under every scheduler \\
T13 & 4-hands piano: Example 9.5 exactly; a third player gives exactly two
  performances \\
T14 & chimera: vibes share 0.9 / 0.5 / 0.1 by articulation; 1 when named \\
\bottomrule
\end{tabularx}
\end{center}

\noindent Measured values for each test are in \ttx{DESIGN.md}.

\end{document}
